% VI32 detailed exercise solutions -- VI/24+ proof-depth standard.

% VI32-DETAILED-EXERCISE-01
\paragraph{Solution VI.32.E1.}
At the origin of
\[
\mathbb A^1_k=\Spec k[t],
\]
the corresponding prime is \((t)\).  Therefore
\[
\mathcal O_{\mathbb A^1,0}=k[t]_{(t)}
\]
and
\[
\dim\mathcal O_{\mathbb A^1,0}
=
\operatorname{ht}(t).
\]
The chain
\[
(0)\subsetneq(t)
\]
has length \(1\), and \(\dim k[t]=1\). Hence
\[
\boxed{\dim\mathcal O_{\mathbb A^1,0}=1}.
\]

% VI32-DETAILED-EXERCISE-02
\paragraph{Solution VI.32.E2.}
At the origin of
\[
\mathbb A^3_k=\Spec k[x,y,z],
\]
the local ring is
\[
k[x,y,z]_{(x,y,z)}.
\]
Its dimension is the height of the maximal ideal.  The chain
\[
(0)\subsetneq(x)\subsetneq(x,y)\subsetneq(x,y,z)
\]
has length \(3\), and the ambient polynomial ring has dimension \(3\). Thus
\[
\boxed{\dim\mathcal O_{\mathbb A^3,0}=3}.
\]

% VI32-DETAILED-EXERCISE-03
\paragraph{Solution VI.32.E3.}
For
\[
\mathfrak p=(x)\subset k[x,y,z],
\]
VI/29 gives
\[
\dim k[x,y,z]_{\mathfrak p}
=
\operatorname{ht}(\mathfrak p).
\]
The prime is principal and nonzero in a domain, so the principal ideal theorem and
\((0)\subsetneq(x)\) give
\[
\operatorname{ht}(x)=1.
\]
Hence
\[
\boxed{\dim\mathcal O_{\Spec k[x,y,z],\mathfrak p}=1}.
\]

% VI32-DETAILED-EXERCISE-04
\paragraph{Solution VI.32.E4.}
Let \(\eta\) be the generic point of an integral variety \(X\).  On a nonempty affine chart
\[
U=\Spec A,
\]
the generic point corresponds to the zero prime because \(A\) is a domain.  Thus
\[
\mathcal O_{X,\eta}\cong A_{(0)}=\Frac(A)=K(X),
\]
which is a field.  Every field has dimension \(0\), so
\[
\boxed{\dim\mathcal O_{X,\eta}=0}.
\]

% VI32-DETAILED-EXERCISE-05
\paragraph{Solution VI.32.E5.}
For a \(5\)-dimensional variety \(X\) and a closed point \(x\), the closed-point local-dimension theorem gives
\[
\dim\mathcal O_{X,x}=\dim X.
\]
Therefore
\[
\boxed{\dim\mathcal O_{X,x}=5}.
\]
The finite-type hypothesis is important: this equality is not true for an arbitrary scheme.

% VI32-DETAILED-EXERCISE-06
\paragraph{Solution VI.32.E6.}
At the origin of \(\mathbb A^4_k\), the maximal ideal is
\[
\mathfrak m=(x_1,x_2,x_3,x_4)
\]
in the local ring.  Modulo \(\mathfrak m^2\), every product of two coordinate functions vanishes.  The four classes
\[
\overline{x}_1,\overline{x}_2,\overline{x}_3,\overline{x}_4
\]
span \(\mathfrak m/\mathfrak m^2\), and no nontrivial \(k\)-linear combination lies in \(\mathfrak m^2\), which
contains only terms of order at least \(2\). Hence
\[
\boxed{\operatorname{edim}_0\mathbb A^4_k=4}.
\]

% VI32-DETAILED-EXERCISE-07
\paragraph{Solution VI.32.E7.}
For
\[
R=k[t]_{(t)},
\qquad
\mathfrak m=(t),
\]
we have
\[
\mathfrak m^2=(t^2).
\]
Thus
\[
\mathfrak m/\mathfrak m^2
=
(t)/(t^2)
\]
is generated by the nonzero class of \(t\).  It is one-dimensional over the residue field \(k\):
\[
\boxed{\dim_k\mathfrak m/\mathfrak m^2=1}.
\]

% VI32-DETAILED-EXERCISE-08
\paragraph{Solution VI.32.E8.}
Let
\[
R=k[\varepsilon]/(\varepsilon^2).
\]
Its unique prime ideal is \((\varepsilon)\), so \(R\) is Artinian and
\[
\dim R=0.
\]
Its maximal ideal is
\[
\mathfrak m=(\varepsilon),
\qquad
\mathfrak m^2=0.
\]
Therefore
\[
\mathfrak m/\mathfrak m^2\cong k
\]
has dimension \(1\). Since
\[
0<1,
\]
the regular-local equality fails. Hence
\[
\boxed{\Spec k[\varepsilon]/(\varepsilon^2)\text{ is not regular}.}
\]

% VI32-DETAILED-EXERCISE-09
\paragraph{Solution VI.32.E9.}
For
\[
X=V(y-x^2)\subseteq\mathbb A^2_k,
\]
the Jacobian row is
\[
(-2x,1).
\]
At the origin this becomes
\[
(0,1),
\]
which has rank \(1\) over every field.  Hence
\[
\dim_kT_0X=2-1=1.
\]
Thus
\[
\boxed{\dim_kT_0X=1}.
\]

% VI32-DETAILED-EXERCISE-10
\paragraph{Solution VI.32.E10.}
For the cusp
\[
f=y^2-x^3,
\]
the gradient is
\[
\nabla f=(-3x^2,2y).
\]
At \((0,0)\) both entries vanish when \(\operatorname{char}k\ne2,3\).  Hence the Jacobian has rank \(0\), and
\[
T_0X=\ker(0:k^2\to k)=k^2.
\]
Therefore
\[
\boxed{\dim_kT_0X=2}.
\]
The local dimension is \(1\), so the cusp point is singular.

% VI32-DETAILED-EXERCISE-11
\paragraph{Solution VI.32.E11.}
For
\[
X=V(xy)\subseteq\mathbb A^2_k,
\]
the Jacobian row is
\[
(y,x).
\]
At the origin it is zero.  Hence the tangent space is all of \(k^2\):
\[
\boxed{\dim_kT_0X=2}.
\]
The local dimension of the crossing is \(1\), so the origin is singular.

% VI32-DETAILED-EXERCISE-12
\paragraph{Solution VI.32.E12.}
The hyperplane
\[
X=V(x)\subseteq\mathbb A^3_k
\]
is isomorphic to \(\mathbb A^2_k\), so its local dimension at the origin is \(2\).  Its Jacobian row is
\[
(1,0,0),
\]
of rank \(1\), so
\[
\dim_kT_0X=3-1=2.
\]
Thus
\[
\boxed{\dim\mathcal O_{X,0}=2=\dim_kT_0X},
\]
and the origin is regular.

% VI32-DETAILED-EXERCISE-13
\paragraph{Solution VI.32.E13.}
For a Noetherian local ring \((R,\mathfrak m,\kappa)\), one always has
\[
\dim R\le\dim_\kappa\mathfrak m/\mathfrak m^2.
\]
The ring is regular precisely when equality holds:
\[
\boxed{
R\text{ regular}
\iff
\dim R=\dim_\kappa\mathfrak m/\mathfrak m^2.
}
\]

% VI32-DETAILED-EXERCISE-14
\paragraph{Solution VI.32.E14.}
Nakayama's lemma shows that
\[
\dim_\kappa\mathfrak m/\mathfrak m^2
\]
equals the minimal number of generators of the maximal ideal \(\mathfrak m\).  Indeed, a set of elements of
\(\mathfrak m\) minimally generates \(\mathfrak m\) exactly when their residue classes form a basis of
\(\mathfrak m/\mathfrak m^2\).

% VI32-DETAILED-EXERCISE-15
\paragraph{Solution VI.32.E15.}
Let \(X=V(f)\subseteq\mathbb A^n_k\) and let \(a\in X(k)\).  If
\[
\nabla f(a)\ne0,
\]
then the Jacobian has rank \(1\), so
\[
\dim_kT_aX=n-1.
\]
A genuine hypersurface has local dimension \(n-1\) at such a point. Therefore tangent and local dimensions agree,
and
\[
\boxed{a\text{ is regular}.}
\]

% VI32-DETAILED-EXERCISE-16
\paragraph{Solution VI.32.E16.}
No.  For a nonzero hypersurface in affine \(n\)-space, the local dimension is \(n-1\).  If all first derivatives
vanish at a rational point \(a\), then the Jacobian has rank \(0\), hence
\[
\dim_kT_aX=n.
\]
Since
\[
n>n-1,
\]
the regular-local equality fails. Therefore such a point is singular.

% VI32-DETAILED-EXERCISE-17
\paragraph{Solution VI.32.E17.}
No.  For a Noetherian local scheme point,
\[
\dim\mathcal O_{X,x}
\le
\operatorname{edim}_xX
=
\dim_{\kappa(x)}T_xX.
\]
Equality is the definition of regularity.  At the cusp origin, for example,
\[
\dim\mathcal O_{C,0}=1
\qquad\text{but}\qquad
\dim T_0C=2.
\]

% VI32-DETAILED-EXERCISE-18
\paragraph{Solution VI.32.E18.}
No.  Reduction preserves prime chains and therefore Krull dimension, but it can change
\(\mathfrak m/\mathfrak m^2\).  For the dual numbers,
\[
X=\Spec k[\varepsilon]/(\varepsilon^2)
\]
has local dimension \(0\) and embedding dimension \(1\), so it is singular.  Its reduction
\[
X_{\mathrm{red}}=\Spec k
\]
has both dimensions \(0\), hence is regular.

% VI32-DETAILED-EXERCISE-19
\paragraph{Solution VI.32.E19.}
The point
\[
x=(t^2+1)\in\Spec\mathbb R[t]
\]
has residue field
\[
\kappa(x)
=
\mathbb R[t]/(t^2+1).
\]
Since \(t^2+1\) is irreducible over \(\mathbb R\), this quotient is
\[
\boxed{\kappa(x)\cong\mathbb C}.
\]

% VI32-DETAILED-EXERCISE-20
\paragraph{Solution VI.32.E20.}
The cotangent space is naturally a vector space over the residue field of the point, not necessarily over the
ground field.  Thus at
\[
x=(t^2+1)\in\Spec\mathbb R[t],
\]
the vector space
\[
\mathfrak m_x/\mathfrak m_x^2
\]
is naturally a vector space over
\[
\boxed{\kappa(x)\cong\mathbb C}.
\]

% VI32-DETAILED-EXERCISE-21
\paragraph{Solution VI.32.E21.}
In
\[
X=\Spec k\sqcup\mathbb A^1_k,
\]
let \(x\) be the point of the first component.  Its local ring is simply
\[
\mathcal O_{X,x}\cong k.
\]
Therefore
\[
\boxed{\dim\mathcal O_{X,x}=0}.
\]

% VI32-DETAILED-EXERCISE-22
\paragraph{Solution VI.32.E22.}
The global dimension of a disjoint union is the maximum of the component dimensions:
\[
\dim X
=
\max\{\dim\Spec k,\dim\mathbb A^1_k\}
=
\max\{0,1\}.
\]
Thus
\[
\boxed{\dim X=1}.
\]
Together with E21 this is the AG8-P3a counterexample.

% VI32-DETAILED-EXERCISE-23
\paragraph{Solution VI.32.E23.}
At the cusp, the local ring is
\[
k[t^2,t^3]_{(t^2,t^3)}.
\]
It has local dimension \(1\) but embedding dimension \(2\), so it is singular.  Normalization replaces it by
\[
k[t]_{(t)},
\]
which has local dimension \(1\) and embedding dimension \(1\).  Thus normalization keeps the dimension but repairs
the nonregular one-dimensional local ring in this example.

% VI32-DETAILED-EXERCISE-24
\paragraph{Solution VI.32.E24.}\mbox{}\par
No.  Both the cusp and the node have two-dimensional Zariski tangent space at the origin as plane curves, but their
lowest nonzero homogeneous terms differ.  For the cusp
\[
y^2-x^3,
\]
the initial form is \(y^2\), a doubled tangent line.  For the node
\[
y^2-x^2(x+1),
\]
the initial degree-two form is
\[
y^2-x^2=(y-x)(y+x),
\]
giving two distinct tangent directions when \(\operatorname{char}k\ne2\).

% VI32-DETAILED-EXERCISE-25
\paragraph{Solution VI.32.E25.}
The preceding chapter,
\[
\boxed{\text{VI/31 — Codimension}},
\]
proved
\[
\operatorname{codim}(Y,X)
=
\dim\mathcal O_{X,\eta_Y}
\]
for an irreducible closed subset \(Y\) with generic point \(\eta_Y\).

% VI32-DETAILED-EXERCISE-26
\paragraph{Solution VI.32.E26.}\mbox{}\par
Among the subjects explicitly deferred beyond VI/32 are Kähler differentials, smooth and étale morphisms,
completions, tangent cones in full generality, Serre's conditions, and divisor theory.  Any two examples suffice;
for instance:
\[
\boxed{\text{Kähler differentials and smooth morphisms}.}
\]
